\documentclass[review,12pt,authoryear]{elsarticle}

%% Packages
\usepackage{lineno}
\usepackage{amsmath,amssymb,amsfonts}
\usepackage{graphicx}
\usepackage{float}
\usepackage{booktabs}
\usepackage{multirow}
\usepackage{algorithm}
\usepackage{algorithmic}
\usepackage{hyperref}
\usepackage{xcolor}
\usepackage{subcaption}
\usepackage{bm}

%% Placeholder command for missing figures/tables
\newcommand{\tbd}[1]{\textcolor{red}{\textbf{[TBD: #1]}}}
\newcommand{\std}[1]{{\scriptsize$\pm$#1}}
\newcommand{\method}{ARR}

\journal{Pattern Recognition}

\begin{document}
\linenumbers

\begin{frontmatter}

\title{TOFM-BP: Transformer Operator Feature Map Backpropagation for Supervised Post-Training}

\author[addr1,addr2,addr3]{sober}
\author[addr1,addr2]{Lei He\corref{cor1}}
\ead{helei2023@tsinghua.edu.cn}

\cortext[cor1]{Corresponding author.}
\address[addr1]{School of Vehicle and Mobility, Tsinghua University, Beijing, China}
\address[addr2]{State Key Laboratory of Intelligent Green Vehicle and Mobility, Tsinghua University, Beijing, China}
\address[addr3]{School of Mathematical Sciences, Beijing University of Posts and Telecommunications, Beijing, China}

\begin{abstract}
Transfer learning with pre-trained Swin Transformers has become a strong paradigm for visual recognition tasks.
However, standard fine-tuning optimizes only the task-specific classification objective and imposes no structural constraint on the self-attention weight matrices.
On small-scale target datasets this unconstrained adaptation tends to drive attention heads into degenerate regimes---overly peaked, nearly rank-deficient, or internally inconsistent with the value representations they aggregate---which correlates with overfitting and degraded generalization.
In this paper, we propose \emph{Attention Reconstruction Regularization} (\method{}), a lightweight auxiliary regularizer that enforces an \emph{algebraic self-consistency} constraint on attention: at each regularized block, the attention output and value matrix should admit a valid row-stochastic reconstruction of the current attention distribution.
Specifically, at selected transformer blocks, we reconstruct the attention weight matrix from the attention output and value matrices via Tikhonov-regularized least squares, project the solution onto the probability simplex using a KKT-based closed-form algorithm, and penalize the discrepancy between the reconstructed and original attention distributions.
\method{} introduces \emph{zero} additional parameters and is applied only during training.
Extensive experiments on seven fine-grained and texture recognition benchmarks with Swin Transformer-Base demonstrate that \method{} consistently improves accuracy on five out of seven datasets, achieving up to $+1.59\%$ on DTD and $+1.51\%$ on FGVC Aircraft, while adding moderate training-only overhead and no inference-time cost.
We further provide comprehensive ablation studies on the loss weight, regularization placement, and warmup scheduling, offering practical guidelines for deploying attention-based regularization in transfer learning.
\end{abstract}

\begin{keyword}
Swin Transformer \sep Attention regularization \sep Transfer learning \sep Fine-tuning \sep Hierarchical transformer
\end{keyword}

\end{frontmatter}

%% ====================================================================
%%  1. INTRODUCTION
%% ====================================================================
\section{Introduction}
\label{sec:intro}

Swin Transformers \citep{liu2021swin}, building on transformer-based image recognition \citep{dosovitskiy2021an}, have established strong results across a broad spectrum of visual recognition tasks.
A prevailing recipe for deploying these models on downstream datasets is \emph{transfer learning}: a backbone pre-trained on a large-scale source domain (\emph{e.g.}, ImageNet-22K) is fine-tuned end-to-end on the target task using a standard cross-entropy loss.
This paradigm has proven remarkably effective, yet it provides no explicit mechanism to constrain the \emph{internal geometry} of the self-attention weight matrices, even though these matrices are the primary site at which overfitting on small target datasets manifests structurally.

Self-attention is the defining computational primitive of Transformers.
In each attention head, a query--key inner product produces a weight matrix $\mathbf{S} \in \mathbb{R}^{N \times N}$ (after softmax normalization) that governs how information from the value tokens is aggregated.
Two structural properties of $\mathbf{S}$ are central to generalization. First, each row of $\mathbf{S}$ lies on the probability simplex $\Delta^{N-1}$, so the output $\mathbf{O} = \mathbf{S}\mathbf{V}$ is a convex combination of value tokens; this constraint is what makes attention interpretable as a routing mechanism rather than an arbitrary linear map. Second, well-functioning attention heads produce distributions that are neither degenerately peaked (which collapses the effective receptive field to a single token) nor nearly uniform (which makes the head an averaging operator and erases token-level discrimination)~\citep{niu2021attention,zhang2024accvit}.
During fine-tuning, however, the attention weights are driven solely by the classification gradient on a small target set, and there is nothing in the loss that prevents the head from drifting into either of these degenerate regimes---or, more subtly, into a regime where the row-stochastic structure of $\mathbf{S}$ becomes \emph{inconsistent} with the linear span of $\mathbf{V}$, in the sense that no valid probability vector over tokens could reproduce the observed output. We argue that this last failure mode---a loss of algebraic self-consistency between $\mathbf{S}$, $\mathbf{V}$, and $\mathbf{O}$---is a useful, previously unexploited signal of attention overfitting.

Several lines of work have sought to regularize the fine-tuning process.
Weight-space regularization methods, such as $L^2$-SP \citep{li2018explicit} and DELTA \citep{li2019delta}, penalize parameter drift from the pre-trained initialization.
Knowledge distillation approaches \citep{hinton2015distilling} use the pre-trained model as a teacher to constrain the fine-tuned model's output distribution.
More recently, attention-level supervision has been explored in the context of self-distillation \citep{zhang2019your} and attention transfer \citep{zagoruyko2017paying}, but these methods typically require a frozen teacher network, introduce additional parameters, or operate at the feature map level rather than directly on attention weight matrices.

In this paper, we operationalize this self-consistency signal directly.
Importantly, our method does \emph{not} compare the fine-tuned attention against a frozen pre-trained copy; it requires no teacher and no reference distribution external to the current forward pass.
Instead, given the attention output $\mathbf{O} = \mathbf{S}\mathbf{V}$ at a selected transformer block, we \emph{reconstruct} the attention weight matrix $\tilde{\mathbf{S}}$ by solving the inverse problem $\mathbf{O} \approx \tilde{\mathbf{S}}\mathbf{V}$ via Tikhonov-regularized least squares, and project the solution onto the probability simplex using a closed-form KKT algorithm \citep{duchi2008efficient}.
The reconstruction loss $\mathcal{L}_{\text{rec}} = \text{MSE}(\tilde{\mathbf{S}},\, \mathbf{S}_{\text{gt}}) \,/\, \mathbb{E}[\mathbf{S}_{\text{gt}}^2]$ measures how faithfully the observed output can be ``explained'' as some valid row-stochastic combination of value tokens that is close to the current attention.
When the head is self-consistent, a near-zero residual is achievable and the loss vanishes; when the head drifts toward a rank-deficient or degenerately peaked configuration, the simplex-constrained reconstruction cannot match the output and the loss grows.
Driving this residual toward zero therefore acts as an \emph{implicit regularizer on the conditioning and smoothness of the attention distribution}---penalizing row-wise peakedness and ill-conditioning in $\mathbf{V}$ that together push $\mathbf{S}$ out of the well-behaved regime---without ever explicitly referencing a target attention map.
This is an orthogonal axis to weight-space anchoring ($L^2$-SP), output-space smoothing (label smoothing, R-Drop), and data-space mixing (CutMix, Mixup), none of which directly constrain the attention geometry.
Figure~\ref{fig:overview} illustrates the overall architecture of the proposed approach.

\begin{figure}[H]
    \centering
    \includegraphics[width=1.05\textwidth]{figures/overview_arr_clean.png}
    \caption{Overview of Attention Reconstruction Regularization (\method{}). (a) The Swin Transformer-Base backbone with four hierarchical stages. (b) Detail of a single transformer block showing the query--key--value projections, dot-product attention with softmax, and the attention output $\mathbf{O} = \mathbf{S}\mathbf{V}$. The ground-truth attention $\mathbf{S}_{\text{gt}}$ is detached (no gradient). (c) The \method{} module: $\mathbf{O}$ and $\mathbf{V}$ are fed into a Tikhonov solve to reconstruct $\tilde{\mathbf{S}}$, which is projected onto the probability simplex via KKT closed-form, and compared against $\mathbf{S}_{\text{gt}}$ using normalized MSE. (d) The classification head produces $\mathcal{L}_{\text{cls}}$, which is combined with $\mathcal{L}_{\text{rec}}$ as $\mathcal{L}_{\text{total}} = (1-\alpha)\mathcal{L}_{\text{cls}} + \alpha\mathcal{L}_{\text{rec}}$.}
    \label{fig:overview}
\end{figure}

We term our approach \emph{Attention Reconstruction Regularization} (\method{}).
It possesses several appealing properties:
\begin{itemize}
    \item \textbf{Zero additional parameters.} \method{} operates entirely on existing intermediate representations ($\mathbf{O}$, $\mathbf{V}$, $\mathbf{S}$) within the attention module and requires no teacher network, projection heads, or auxiliary modules.
    \item \textbf{Training-time only.} The reconstruction loss is computed only during training and adds no cost at inference.
    \item \textbf{Flexible placement.} The regularization can be applied to any subset of transformer blocks, enabling fine-grained control over the trade-off between regularization strength and computational overhead.
    \item \textbf{Complementary to existing techniques.} \method{} regularizes a different aspect of the model (attention structure) than weight decay, dropout, or output-level distillation, and can in principle be combined with these methods.
\end{itemize}

We evaluate \method{} by fine-tuning a Swin Transformer-Base (pre-trained on ImageNet-22K) on seven diverse benchmarks spanning fine-grained recognition (CUB-200-2011, Stanford Cars, FGVC Aircraft), texture classification (DTD), pet breed classification (Oxford-IIIT Pet), flower classification (Flowers-102), and dog breed classification (Stanford Dogs).
\method{} yields consistent improvements on five of the seven datasets, with gains of up to $+1.59\%$ on DTD and $+1.51\%$ on FGVC Aircraft under the best configuration, while the two datasets where it does not help (Flowers-102 at 99.4\% accuracy, Stanford Dogs) present informative boundary cases that we analyze in detail.
Comprehensive ablation studies examine the sensitivity to the loss weight $\alpha$, the choice of which transformer blocks to regularize, and the effect of a cosine warmup schedule for $\alpha$.

The contributions of this paper are summarized as follows:
\begin{enumerate}
    \item We formulate a novel attention-level self-consistency regularization for Swin Transformer fine-tuning, based on reconstructing attention weights from the output--value linear system via Tikhonov regularization and simplex projection.
    \item We conduct extensive experiments on seven benchmarks and provide practical guidelines on loss weight selection ($\alpha = 0.1$), regularization placement (penultimate stage preferred), and warmup scheduling (beneficial for small datasets with fewer than 4{,}000 training samples).
    \item We offer a frank analysis of failure modes---accuracy saturation and intra-class variability---that delineates the applicability boundary of attention-based regularization in transfer learning.
\end{enumerate}

\begin{figure}[H]
    \centering
    \includegraphics[width=1.05\textwidth]{figures/pipeline.png}
    \caption{Simplified pipeline of the \method{} reconstruction loss. From a selected Swin Transformer block, the attention output $\mathbf{O}$ and value matrix $\mathbf{V}$ are used to solve for $\tilde{\mathbf{S}}$ via Tikhonov regularization, projected onto the probability simplex $\Delta_{N-1}$ using KKT closed-form, and compared against the detached ground-truth $\mathbf{S}_{\text{gt}}$ via normalized MSE. The reconstruction loss $\mathcal{L}_{\text{rec}}$ is combined with the classification loss $\mathcal{L}_{\text{cls}}$ to form the total training objective.}
    \label{fig:pipeline}
\end{figure}

The remainder of this paper is organized as follows.
Section~\ref{sec:related} reviews related work on Swin Transformers, transfer learning regularization, and attention supervision.
Section~\ref{sec:method} details the proposed \method{} method.
Section~\ref{sec:experiments} presents the experimental setup and results.
Section~\ref{sec:discussion} discusses the findings, and Section~\ref{sec:conclusion} concludes the paper.

%% ====================================================================
%%  2. RELATED WORK
%% ====================================================================
\section{Related Work}
\label{sec:related}

We situate our work at the intersection of three research threads: Swin Transformer architectures, regularization strategies for transfer learning, and attention-level supervision.

%% --------------------------------------------------------------------
\subsection{Swin Transformers}
\label{sec:related:vit}

\citet{dosovitskiy2021an} introduced the Vision Transformer (ViT), which applies a standard Transformer encoder \citep{vaswani2017attention} to sequences of image patches, achieving competitive performance with convolutional networks when pre-trained on large-scale datasets.
The Swin Transformer \citep{liu2021swin} adapts this paradigm into a hierarchical architecture with shifted window-based self-attention, reducing the quadratic complexity of global attention to linear in the image size while maintaining a multi-scale feature pyramid.
This makes Swin Transformer especially suitable for dense visual representations and transfer learning under limited target data.
Related hierarchical and efficient transformer variants include PVT \citep{wang2021pyramid}, TrpViT \citep{li2024trpvit}, and local-inductive-bias transformer models \citep{akkaya2024enhancing}, each proposing different spatial reduction, local bias, or attention partitioning schemes.
Our work builds on the Swin Transformer but the proposed regularization is architecture-agnostic: it applies to any transformer block that computes $\mathbf{O} = \mathbf{S}\mathbf{V}$.

%% --------------------------------------------------------------------
\subsection{Regularization for Transfer Learning}
\label{sec:related:transfer}

Transfer learning via fine-tuning a pre-trained backbone is the de facto standard for visual recognition with limited target data \citep{zhuang2020comprehensive,akkaya2024enhancing}.
A central challenge is to transfer pre-trained knowledge effectively while adapting to the target distribution, especially when the target dataset is small and overfitting is likely.

\textbf{Parameter-space regularization.}
$L^2$-SP \citep{li2018explicit} penalizes the $L^2$ distance between fine-tuned and pre-trained weights, acting as a ``soft anchor'' to the initialization.
DELTA \citep{li2019delta} extends this idea by regularizing in the feature space through attention maps derived from intermediate activations.
BSS \citep{chen2019catastrophic} suppresses small singular values of the classifier weight matrix to mitigate negative transfer.
These methods operate on weights or features but do not directly constrain the attention mechanism.

\textbf{Parameter-efficient fine-tuning.}
An alternative line of work freezes most pre-trained parameters and inserts small trainable modules.
Adapters \citep{houlsby2019parameter}, LoRA \citep{hu2022lora}, and visual prompt tuning (VPT) \citep{jia2022visual} add fewer than 1\% extra parameters while achieving competitive accuracy.
These approaches implicitly preserve pre-trained representations by limiting the degrees of freedom, but they sacrifice full end-to-end adaptation and do not exploit the internal structure of attention.

\textbf{Output-space regularization.}
Knowledge distillation \citep{hinton2015distilling} uses the softened output probabilities of a teacher network to guide the student.
When the teacher is the pre-trained model itself, this becomes a form of self-distillation that regularizes the fine-tuning process \citep{mobahi2020self}.
R-Drop \citep{liang2021r} regularizes output consistency between two forward passes with different dropout masks.
Label smoothing \citep{szegedy2016rethinking} and data augmentation strategies such as CutMix \citep{yun2019cutmix} and Mixup \citep{zhang2018mixup} provide complementary regularization.
Our method differs from all of the above in that it regularizes at the \emph{attention weight} level---a structural layer of the model that is orthogonal to parameter drift, output logits, or data perturbation.

%% --------------------------------------------------------------------
\subsection{Attention Supervision and Distillation}
\label{sec:related:attn}

The idea of supervising or transferring attention distributions has been explored in multiple contexts.

\textbf{Attention transfer.}
\citet{zagoruyko2017paying} proposed transferring spatial attention maps (channel-wise pooled activation norms) from a teacher to a student CNN.
This was later extended to self-attention in transformer-based recognition methods such as \citet{yu2023mixvit}, which emphasize the interaction between discriminative parts and compact transformer representations.
These methods require a separate, frozen teacher network, which doubles the memory footprint during training.

\textbf{Self-distillation.}
\citet{zhang2019your} showed that a network can distill knowledge from its own deeper layers to shallower ones, improving both accuracy and calibration.
Granularity-aware distillation has been used to stabilize fine-grained visual recognition across feature levels \citep{ke2023granularity}.
Token-based self-supervised pre-training has also been revisited for transformer backbones, showing that reconstruction-style proxy tasks can improve downstream transfer \citep{tian2025beyond}.
While self-distillation shares our spirit of exploiting a model's own representations, it typically operates at the logit or feature level and requires architectural modifications (\emph{e.g.}, auxiliary classifiers at intermediate layers).

\textbf{Attention regularization.}
More closely related to our work, several studies have imposed explicit constraints on attention matrices.
\citet{niu2021attention} modeled attention shifts to iteratively discover discriminative regions for fine-grained recognition.
\citet{zhang2022sequential} promoted spatially diversified representations through sequential attention constraints.
Recent transformer-based fine-grained recognition methods further use token selection, attention aggregation, or token mixing to emphasize discriminative regions \citep{yu2023mixvit,wang2023aatrans,zhang2024accvit}.
Self-supervised part erasing and random mask covariance learning provide related ways to strengthen part-level representations from limited data \citep{yu2021maskcov,yu2022spare}.
These methods modify the representation pipeline, use additional branches, or introduce task-specific attention modules.
In contrast, \method{} does not alter the model's forward pass at all; it adds a purely auxiliary loss that is derived from the \emph{inverse problem} of reconstructing attention weights from outputs.

\textbf{Inverse problems in neural networks.}
The idea of reconstructing intermediate representations from downstream outputs has connections to invertible networks \citep{jacobsen2018revnet} and feature reconstruction in autoencoders.
Our formulation---solving $\mathbf{O} \approx \tilde{\mathbf{S}}\mathbf{V}$ for $\tilde{\mathbf{S}}$---is a constrained linear inverse problem.
The Tikhonov regularization \citep{tikhonov1977solutions} ensures numerical stability, and the simplex projection via KKT conditions follows the efficient algorithm of \citet{duchi2008efficient}.
To the best of our knowledge, this is the first work to formulate attention regularization as an inverse reconstruction problem and to apply it to Swin Transformer fine-tuning.

%% ====================================================================
%%  3. METHOD
%% ====================================================================
\section{Method}
\label{sec:method}

This section first reviews the self-attention mechanism in the Swin Transformer (Section~\ref{sec:method:prelim}), then presents the attention reconstruction formulation (Section~\ref{sec:method:recon}), the simplex projection algorithm (Section~\ref{sec:method:simplex}), the normalized reconstruction loss (Section~\ref{sec:method:loss}), and the overall training objective with warmup scheduling (Section~\ref{sec:method:total}).

%% --------------------------------------------------------------------
\subsection{Preliminaries: Window-based Self-Attention}
\label{sec:method:prelim}

The Swin Transformer \citep{liu2021swin} partitions each feature map into non-overlapping local windows of size $M \times M$ (default $M=7$) and computes self-attention independently within each window.
Let $N = M^2$ denote the number of tokens per window.
For each attention head $h$ in a given transformer block, the input tokens $\mathbf{X} \in \mathbb{R}^{N \times C}$ are projected into queries, keys, and values:
\begin{equation}
    \mathbf{Q} = \mathbf{X}\mathbf{W}_Q, \quad
    \mathbf{K} = \mathbf{X}\mathbf{W}_K, \quad
    \mathbf{V} = \mathbf{X}\mathbf{W}_V,
    \label{eq:qkv}
\end{equation}
where $\mathbf{Q}, \mathbf{K}, \mathbf{V} \in \mathbb{R}^{N \times d}$ and $d = C / n_H$ is the per-head dimension.
The attention weight matrix and output are computed as:
\begin{equation}
    \mathbf{S} = \mathrm{softmax}\!\left(\frac{\mathbf{Q}\mathbf{K}^\top}{\sqrt{d}} + \mathbf{B}\right), \quad
    \mathbf{O} = \mathbf{S}\mathbf{V},
    \label{eq:attn}
\end{equation}
where $\mathbf{B}$ is the learnable relative position bias and $\mathbf{S} \in \mathbb{R}^{N \times N}$ satisfies $\mathbf{S} \geq 0$, $\mathbf{S}\mathbf{1} = \mathbf{1}$ (each row lies on the probability simplex $\Delta^{N-1}$).

In the Swin-Base configuration, the four stages have depths $(2, 2, 18, 2)$ with $(4, 8, 16, 32)$ attention heads and embedding dimensions $(128, 256, 512, 1024)$, respectively.
The total number of parameters is 86.8M.

%% --------------------------------------------------------------------
\subsection{Attention Weight Reconstruction}
\label{sec:method:recon}

The key observation underlying \method{} is that the relationship $\mathbf{O} = \mathbf{S}\mathbf{V}$ defines a linear system in the unknown $\mathbf{S}$, given the observed output $\mathbf{O}$ and value matrix $\mathbf{V}$.
In principle, if $\mathbf{V}$ has full row rank, $\mathbf{S}$ can be recovered exactly.
In practice, $\mathbf{V} \in \mathbb{R}^{N \times d}$ with $N > d$ (\emph{e.g.}, $N = 49$, $d = 32$), so the system is underdetermined per row of $\mathbf{S}$ and ill-conditioned.
We therefore adopt Tikhonov-regularized least squares:
\begin{equation}
    \tilde{\mathbf{S}} = \arg\min_{\mathbf{S}'} \left\| \mathbf{O} - \mathbf{S}'\mathbf{V} \right\|_F^2 + \lambda \left\| \mathbf{S}' \right\|_F^2,
    \label{eq:tikhonov}
\end{equation}
which admits the closed-form solution:
\begin{equation}
    \tilde{\mathbf{S}} = \mathbf{O}\mathbf{V}^\top \left(\mathbf{V}\mathbf{V}^\top + (\lambda + \varepsilon)\mathbf{I}\right)^{-1},
    \label{eq:solve_S}
\end{equation}
where $\lambda = 0.05$ is the Tikhonov regularization coefficient and $\varepsilon = 10^{-3}$ is a numerical stability term.
The matrix inversion operates on $N \times N$ matrices ($49 \times 49$ for window size 7), which is computationally inexpensive.
In practice, we use \texttt{torch.linalg.solve} instead of explicit inversion for better numerical properties.

\textbf{Gradient flow.}
A critical design choice is that $\tilde{\mathbf{S}}$ is computed from $\mathbf{O}$ and $\mathbf{V}$, both of which are differentiable functions of the model parameters.
The ground-truth target $\mathbf{S}_{\text{gt}} = \mathbf{S}.\text{detach}()$ is treated as a fixed supervision signal (no gradient flows through $\mathbf{S}_{\text{gt}}$).
This means the reconstruction loss encourages the value and output representations to be mutually consistent in a way that aligns with the current attention distribution, without directly constraining the query--key interaction.

%% --------------------------------------------------------------------
\subsection{Simplex Projection via KKT Conditions}
\label{sec:method:simplex}

The unconstrained solution $\tilde{\mathbf{S}}$ from Eq.~\eqref{eq:solve_S} does not in general satisfy the probability simplex constraints ($\tilde{\mathbf{S}} \geq 0$, row sums equal to 1).
To enforce these constraints, we project each row $\tilde{\mathbf{s}}_i \in \mathbb{R}^N$ onto the simplex $\Delta^{N-1}$:
\begin{equation}
    \mathbf{s}_i^* = \arg\min_{\mathbf{s} \in \Delta^{N-1}} \left\| \mathbf{s} - \tilde{\mathbf{s}}_i \right\|_2^2.
    \label{eq:simplex_proj}
\end{equation}
This is a convex quadratic program with a well-known closed-form solution based on the KKT conditions \citep{duchi2008efficient}.
The KKT stationarity condition yields $s_j^* = \max(\tilde{s}_{i,j} - \theta, \, 0)$, where the threshold $\theta$ is determined by the equality constraint $\sum_j s_j^* = 1$.

\textbf{KKT derivation for simplex projection.}
For a row vector $\hat{\mathbf{s}} = \tilde{\mathbf{s}}_i$, the Euclidean projection solves
\begin{equation}
    \min_{\mathbf{s}} \frac{1}{2}\|\mathbf{s} - \hat{\mathbf{s}}\|_2^2
    \quad \text{s.t.} \quad
    \sum_j s_j = 1,\; s_j \geq 0.
    \label{eq:simplex_kkt_obj}
\end{equation}
The Lagrangian is
\begin{equation}
    \mathcal{L}(\mathbf{s}, \theta, \bm{\nu})
    = \frac{1}{2}\|\mathbf{s} - \hat{\mathbf{s}}\|_2^2
    + \theta\left(\sum_j s_j - 1\right)
    - \sum_j \nu_j s_j,
    \label{eq:simplex_lagrangian}
\end{equation}
where $\theta$ enforces the equality constraint and $\nu_j \geq 0$ enforces non-negativity.
The KKT conditions are
\begin{align}
    s_j - \hat{s}_j + \theta - \nu_j &= 0, \label{eq:simplex_stationarity} \\
    s_j &\geq 0,\quad \nu_j \geq 0,\quad \nu_j s_j = 0, \label{eq:simplex_complementarity} \\
    \sum_j s_j &= 1. \label{eq:simplex_sum}
\end{align}
If $s_j > 0$, complementary slackness gives $\nu_j = 0$ and hence $s_j = \hat{s}_j - \theta$; if $s_j = 0$, the same stationarity condition implies $\hat{s}_j - \theta \leq 0$.
Therefore,
\begin{equation}
    s_j^* = \max(\hat{s}_j - \theta, 0),
    \quad
    \sum_j \max(\hat{s}_j - \theta, 0) = 1.
    \label{eq:simplex_solution}
\end{equation}

The complete algorithm proceeds as follows for each row:
\begin{enumerate}
    \item Sort the entries in descending order: $u_1 \geq u_2 \geq \cdots \geq u_N$.
    \item Compute the cumulative sums $C_k = \sum_{j=1}^{k} u_j$.
    \item Find $\rho = \max\left\{k \,:\, u_k > \frac{C_k - 1}{k}\right\}$.
    \item Set $\theta = \frac{C_\rho - 1}{\rho}$.
    \item Return $s_j^* = \max(\tilde{s}_{i,j} - \theta, \, 0)$.
\end{enumerate}
The per-row complexity is $O(N \log N)$ due to sorting.
Since $N = 49$ in the Swin Transformer with window size 7, this projection is negligible compared to the attention computation itself.
The projection is fully differentiable (the $\max$ operation is piecewise linear), so gradients flow through it during backpropagation.

\textbf{Optional extension: margin-constrained projection.}
The simplex projection above is the only KKT projection used in the proposed \method{} loss and in all experiments in this paper.
As a possible advanced extension, the same projection viewpoint can be applied to a different problem: projecting a score vector onto a classification-margin feasible set.
This extension is not part of the current \method{} training objective, but we include it to clarify how the inverse-reconstruction idea could be coupled with stronger class-level constraints in future work.

Let $\mathbf{q}^*$ denote the optimal projected score vector, $\hat{\mathbf{q}}$ the input score vector, and $j^*$ the target class index.
The margin-constrained projection solves the following $L_2$-norm optimization problem:
\begin{equation}
    \min_{\mathbf{q}^*} \frac{1}{2}\|\mathbf{q}^* - \hat{\mathbf{q}}\|_2^2 \quad \text{s.t.} \quad q_{j^*}^* - q_i^* \geq \Delta, \; \forall i \neq j^*.
    \label{eq:kkt_obj}
\end{equation}
This is a convex quadratic program satisfying the Slater condition, so the optimal solution $\mathbf{q}^*$ exists and is unique.

\textbf{Lagrangian and KKT conditions.}
Introducing Lagrange multipliers $\mu_i \geq 0$ for $i \neq j^*$, the Lagrangian is:
\begin{equation}
    \mathcal{L}(\mathbf{q}^*, \bm{\mu}) = \frac{1}{2}\|\mathbf{q}^* - \hat{\mathbf{q}}\|_2^2 - \sum_{i \neq j^*} \mu_i(q_{j^*}^* - q_i^* - \Delta).
    \label{eq:kkt_lagrangian}
\end{equation}
The KKT conditions are:
\begin{align}
    \frac{\partial \mathcal{L}}{\partial q_i^*} &= (q_i^* - \hat{q}_i) + \mu_i = 0, \quad i \neq j^*, \label{eq:kkt_stat_i} \\
    \frac{\partial \mathcal{L}}{\partial q_{j^*}^*} &= (q_{j^*}^* - \hat{q}_{j^*}) - \sum_{i \neq j^*} \mu_i = 0, \label{eq:kkt_stat_j} \\
    \mu_i(q_{j^*}^* - q_i^* - \Delta) &= 0, \quad \forall i \neq j^*. \label{eq:kkt_comp}
\end{align}

\textbf{Eliminating primal variables.}
From Eqs.~\eqref{eq:kkt_stat_i}--\eqref{eq:kkt_stat_j}, we obtain:
\begin{align}
    q_i^* &= \hat{q}_i - \mu_i, \quad i \neq j^*, \label{eq:kkt_qi} \\
    q_{j^*}^* &= \hat{q}_{j^*} + \sum_{i \neq j^*} \mu_i. \label{eq:kkt_qj}
\end{align}
Substituting into the constraint yields:
\begin{equation}
    \hat{q}_{j^*} - \hat{q}_i + \sum_{k \neq j^*} \mu_k + \mu_i \geq \Delta.
    \label{eq:kkt_substituted}
\end{equation}

\textbf{Active set partition.}
Define the active set:
\begin{equation}
    I_3 = \{i \neq j^* \mid q_{j^*}^* - q_i^* = \Delta\}.
    \label{eq:active_set}
\end{equation}
The remaining index sets are:
\begin{align}
    I_1 &= \{i \mid \hat{q}_{j^*} - \hat{q}_i \geq \Delta\}, \label{eq:I1} \\
    I_2 &= \{i \mid \hat{q}_{j^*} - \hat{q}_i < \Delta\}. \label{eq:I2}
\end{align}
For $i \in I_1$, the constraint is inactive and $\mu_i = 0$; for $i \in I_3$, the constraint is active and $\mu_i > 0$.

\textbf{Closed-form solution on the active set.}
For any $i \in I_3$, the equality constraint gives:
\begin{equation}
    \sum_{k \in I_3} \mu_k + \mu_i = \Delta - (\hat{q}_{j^*} - \hat{q}_i).
    \label{eq:active_eq}
\end{equation}
Summing over $i \in I_3$:
\begin{equation}
    (|I_3| + 1)\sum_{k \in I_3} \mu_k = \sum_{k \in I_3}\bigl(\Delta - (\hat{q}_{j^*} - \hat{q}_k)\bigr).
    \label{eq:sum_mu}
\end{equation}
Therefore:
\begin{equation}
    \sum_{k \in I_3} \mu_k = \frac{\sum_{k \in I_3}\bigl(\Delta - (\hat{q}_{j^*} - \hat{q}_k)\bigr)}{|I_3| + 1}.
    \label{eq:sum_mu_solved}
\end{equation}
The individual multipliers are:
\begin{equation}
    \mu_i = \frac{\Delta - (\hat{q}_{j^*} - \hat{q}_i)}{|I_3| + 1} - \frac{1}{|I_3| + 1}\sum_{k \in I_3}\bigl(\Delta - (\hat{q}_{j^*} - \hat{q}_k)\bigr), \quad i \in I_3.
    \label{eq:mu_i_solved}
\end{equation}

Figure~\ref{fig:framework_inverse} provides a visual summary of this optional extended pipeline.
Panels (b)--(d), including inverse softmax, query reconstruction, and input-space margin projection, are not used in the main \method{} loss or in the experiments reported in this paper; they illustrate a possible direction for future margin-constrained variants.

\begin{figure}[H]
    \centering
    \includegraphics[width=\textwidth]{figures/framework_inverse_tikz.pdf}
    \caption{Optional extension beyond the main \method{} pipeline. The method evaluated in this paper uses only attention reconstruction followed by row-wise simplex projection of $\mathbf{S}$, as described in Sections~\ref{sec:method:recon}--\ref{sec:method:simplex}. Panels (b)--(d) illustrate a possible advanced variant: inverse softmax from $\mathbf{S}$ to log-space $\mathbf{Z}$, query reconstruction from $\mathbf{Z}$, and input-space projection via KKT conditions to satisfy classification margin constraints. This margin-projection branch is not used in the reported experiments.}
    \label{fig:framework_inverse}
\end{figure}

%% --------------------------------------------------------------------
\subsection{Normalized Reconstruction Loss}
\label{sec:method:loss}

After projection, we obtain $\tilde{\mathbf{S}} \in \Delta^{N-1}$ (per row) and compare it against the ground-truth attention weights $\mathbf{S}_{\text{gt}} = \mathbf{S}.\text{detach}()$.
The raw MSE between the two matrices can vary by orders of magnitude across layers and heads due to differences in the sharpness of the attention distribution (\emph{e.g.}, a nearly uniform distribution has $\mathbf{S}_{ij} \approx 1/N$, while a peaked distribution has entries close to 1).
To normalize across these variations, we adopt a \emph{relative MSE}:
\begin{equation}
    \mathcal{L}_{\text{rec}} = \frac{\left\| \tilde{\mathbf{S}} - \mathbf{S}_{\text{gt}} \right\|_F^2 \,/\, n}
    {\mathbb{E}\!\left[\mathbf{S}_{\text{gt}}^2\right]},
    \label{eq:rec_loss}
\end{equation}
where $n$ is the total number of elements and $\mathbb{E}[\mathbf{S}_{\text{gt}}^2] = \|\mathbf{S}_{\text{gt}}\|_F^2 / n$ is the mean squared value of the ground-truth attention.
This normalization ensures that $\mathcal{L}_{\text{rec}}$ is scale-invariant: a perfect reconstruction yields $\mathcal{L}_{\text{rec}} = 0$ regardless of the attention distribution shape, and a constant fractional error produces a consistent loss magnitude across layers.

When multiple transformer blocks are selected for regularization (see Section~\ref{sec:method:total}), the per-block losses are averaged:
\begin{equation}
    \bar{\mathcal{L}}_{\text{rec}} = \frac{1}{|\mathcal{B}|} \sum_{b \in \mathcal{B}} \mathcal{L}_{\text{rec}}^{(b)},
    \label{eq:rec_loss_avg}
\end{equation}
where $\mathcal{B}$ denotes the set of selected blocks.

%% --------------------------------------------------------------------
\subsection{Total Loss and Warmup Scheduling}
\label{sec:method:total}

The total training loss combines the standard cross-entropy classification loss with the reconstruction regularization:
\begin{equation}
    \mathcal{L}_{\text{total}} = (1 - \alpha)\,\mathcal{L}_{\text{cls}} + \alpha\,\bar{\mathcal{L}}_{\text{rec}},
    \label{eq:total_loss}
\end{equation}
where $\alpha \in [0, 1)$ controls the regularization strength.
Our experiments identify $\alpha = 0.1$ as a robust default across datasets (Section~\ref{sec:exp:alpha}).

\textbf{Block selection (\texttt{recon\_blocks}).}
The reconstruction loss can be applied to any subset of transformer blocks.
We consider five configurations:
\begin{itemize}
    \item \texttt{last\_last}: the last block of the final stage (Stage~4).
    \item \texttt{second\_last}: the last block of the penultimate stage (Stage~3).
    \item \texttt{last\_all}: all blocks of the final stage.
    \item \texttt{all\_last}: the last block of every stage (Stages~1--4).
    \item \texttt{last\_second}: the second-to-last block of the final stage.
\end{itemize}
Using a single block (\texttt{last\_last} or \texttt{second\_last}) keeps the overhead lower than multi-block configurations, which provide broader regularization coverage at the cost of additional memory and computation.

\textbf{Warmup scheduling.}
On small target datasets, the classification loss has not yet converged in early epochs, and a large reconstruction weight may interfere with feature adaptation.
To address this, we introduce a cosine warmup schedule that ramps $\alpha$ from 0 to $\alpha_{\max}$ over the course of training:
\begin{equation}
    \alpha(t) = \frac{\alpha_{\max}}{2}\left(1 - \cos\!\left(\frac{\pi\, t}{T}\right)\right),
    \label{eq:warmup}
\end{equation}
where $t$ is the current epoch (0-indexed) and $T$ is the total number of epochs.
This ensures $\alpha(0) = 0$ (pure classification in the first epoch) and $\alpha(T) \approx \alpha_{\max}$ (full regularization strength at the end).
Warmup is especially beneficial for datasets with fewer than 4{,}000 training samples (Section~\ref{sec:exp:warmup}).

\textbf{Computational overhead.}
The dominant additional cost is the matrix solve in Eq.~\eqref{eq:solve_S}, which involves a $49 \times 49$ linear system per head per window.
This computation is performed in \texttt{float32} (even under mixed-precision training) to ensure numerical stability.
In practice, on 4$\times$ NVIDIA A100 GPUs, the peak GPU memory increases from approximately 11\,GB to 14.7\,GB per GPU ($+$3.5\,GB), and training throughput decreases by 15--20\%.
The overhead is entirely at training time; inference is unaffected.

%% ====================================================================
%%  4. EXPERIMENTS (original)
%% ====================================================================
\section{Experiments}
\label{sec:experiments}

%% --------------------------------------------------------------------
\subsection{Experimental Setup}
\label{sec:exp:setup}

\textbf{Backbone.}
We use the Swin Transformer-Base \citep{liu2021swin} pre-trained on ImageNet-22K (86.8M parameters, patch size 4, window size 7, input resolution $224 \times 224$).
The pre-trained classification head is discarded and replaced with a new linear head matching the number of target classes.

\textbf{Datasets.}
We evaluate on seven benchmarks spanning fine-grained recognition, texture classification, and general visual recognition.
Table~\ref{tab:datasets} summarizes the dataset statistics.

\begin{table}[t]
\centering
\caption{Dataset statistics. ``Train'' includes all images used for training (official train + val splits merged where applicable).}
\label{tab:datasets}
\small
\begin{tabular}{lcccc}
\toprule
Dataset & Classes & Train & Test & Domain \\
\midrule
Oxford-IIIT Pet \citep{parkhi2012cats}  & 37  & 3{,}680  & 3{,}669  & Pet breeds \\
DTD \citep{cimpoi2014describing}         & 47  & 3{,}760  & 1{,}880  & Textures \\
CUB-200-2011 \citep{wah2011caltech}      & 200 & 5{,}994  & 5{,}794  & Bird species \\
Stanford Cars \citep{krause2013collecting} & 196 & 8{,}144  & 8{,}041  & Car models \\
FGVC Aircraft \citep{maji2013fine}        & 100 & 6{,}667  & 3{,}333  & Aircraft variants \\
Flowers-102 \citep{nilsback2008automated} & 102 & 1{,}020  & 6{,}149  & Flower species \\
Stanford Dogs \citep{khosla2011novel}     & 120 & 12{,}000 & 8{,}580  & Dog breeds \\
\bottomrule
\end{tabular}
\end{table}

\textbf{Training details.}
All models are trained with the AdamW optimizer \citep{loshchilov2019decoupled} with learning rate $10^{-4}$ and weight decay 0.05.
We use cosine annealing learning rate scheduling (no warmup epochs for the learning rate) for 30 epochs.
The batch size is 64 per GPU on 2 NVIDIA A100-40GB GPUs with Distributed Data Parallel (DDP) training via \texttt{torchrun}.
Mixed-precision training (AMP) is enabled, except that the reconstruction loss computation (Eq.~\eqref{eq:solve_S}) is forced to \texttt{float32} for numerical stability.
Data augmentation includes \texttt{RandomResizedCrop}(224), horizontal flip, \texttt{ColorJitter}, and random rotation ($\pm 20^\circ$) for training; \texttt{Resize}(256) followed by \texttt{CenterCrop}(224) for evaluation.
All experiments are conducted under identical initialization and training conditions for reproducibility.
All results are averaged over 3 independent runs with different random seeds.
The reported uncertainty denotes standard deviation.
We additionally perform paired t-tests between baseline and \method{}.

\textbf{Default \method{} configuration.}
Unless stated otherwise, the reconstruction loss uses $\alpha = 0.1$, \texttt{recon\_blocks} = \texttt{second\_last} (Stage~3 last block), Tikhonov $\lambda = 0.05$, $\varepsilon = 10^{-3}$, KKT simplex projection enabled, and normalized MSE.
No warmup scheduling is applied by default.

%% --------------------------------------------------------------------
\subsection{Main Results}
\label{sec:exp:main}

Table~\ref{tab:main} reports the classification accuracy of the baseline (standard fine-tuning) and \method{} with the default configuration (\texttt{second\_last}, $\alpha = 0.1$) across all seven datasets.

\begin{table}[t]
\centering
\caption{Main results: classification accuracy (\%) with the default \method{} configuration (\texttt{second\_last}, $\alpha = 0.1$). Bold indicates improvement over baseline. Best result per dataset including all configurations is shown in the rightmost column.}
\label{tab:main}
\footnotesize
\setlength{\tabcolsep}{3pt}
\begin{tabular}{lccccc}
\toprule
Dataset & Baseline & \method{} (default) & $\Delta$ & Best \method{} & Best $\Delta$ \\
\midrule
Oxford-IIIT Pet  & 93.37\std{0.04} & \textbf{93.92}\std{0.05} & +0.55\std{0.03} & 94.14\std{0.03}$^{\dagger}$  & +0.77\std{0.04} \\
DTD              & 81.17\std{0.07} & \textbf{82.45}\std{0.04} & +1.28\std{0.05} & 82.76\std{0.06}$^{\ddagger}$ & +1.59\std{0.07} \\
CUB-200-2011     & 89.63\std{0.03} & \textbf{89.82}\std{0.05} & +0.19\std{0.04} & 89.82\std{0.04}              & +0.19\std{0.04} \\
Stanford Cars    & 89.95\std{0.06} & \textbf{90.82}\std{0.03} & +0.87\std{0.05} & 90.82\std{0.05}              & +0.87\std{0.05} \\
FGVC Aircraft    & 73.42\std{0.08} & \textbf{74.51}\std{0.06} & +1.09\std{0.06} & 74.93\std{0.04}$^{\S}$       & +1.51\std{0.06} \\
Flowers-102      & 99.35\std{0.02} & 99.22\std{0.03}          & $-$0.13\std{0.03} & 99.35\std{0.02}            & 0.00\std{0.02} \\
Stanford Dogs    & 86.48\std{0.05} & 86.09\std{0.07}          & $-$0.39\std{0.05} & 86.48\std{0.04}            & 0.00\std{0.03} \\
\bottomrule
\multicolumn{6}{l}{\footnotesize $^{\dagger}$\texttt{last\_all}, $\alpha=0.1$. \quad $^{\ddagger}$\texttt{second\_last}, $\alpha$ warmup. \quad $^{\S}$\texttt{last\_last}, $\alpha=0.1$.}
\end{tabular}
\end{table}

\method{} improves accuracy on five of seven datasets under the default configuration.
The largest gains occur on DTD ($+$1.28\%) and FGVC Aircraft ($+$1.09\%), both of which are texture- or fine-grained-recognition tasks with moderate training set sizes.
Stanford Cars also benefits substantially ($+$0.87\%).
On Flowers-102, accuracy is already near the ceiling (99.35\%), leaving little room for improvement; \method{}'s marginal interference ($-$0.13\%) is practically negligible.
Stanford Dogs is the only dataset with a clear negative effect ($-$0.39\%), which we analyze in Section~\ref{sec:discussion}.

When dataset-specific tuning of \texttt{recon\_blocks} and warmup is allowed (``Best \method{}'' column), the gains increase further: $+$1.59\% on DTD (with warmup) and $+$1.51\% on Aircraft (\texttt{last\_last}).

To provide visual intuition for these gains, Figure~\ref{fig:attention} compares Grad-CAM saliency maps between the baseline and \method{} on FGVC Aircraft test images where the baseline misclassifies but \method{} predicts correctly.
\method{} tends to consolidate saliency onto the aircraft body and reduce mass on background structures, consistent with the design intent of the reconstruction objective.

\begin{figure}[H]
    \centering
    \includegraphics[width=0.95\textwidth]{figures/attention_comparison.png}
    \caption{Grad-CAM saliency comparison on FGVC Aircraft test images, computed at the last block of Swin stage~3 ($14 \times 14$). Each row is an example where the baseline misclassifies but \method{} predicts correctly. \method{} tends to pull residual saliency away from background structures and concentrate it more tightly on the aircraft body.}
    \label{fig:attention}
\end{figure}

%% --------------------------------------------------------------------
\subsection{Ablation: Loss Weight $\alpha$}
\label{sec:exp:alpha}

We systematically vary $\alpha$ on the DTD dataset using \texttt{last\_last} placement.
Table~\ref{tab:alpha} reports the best accuracy achieved over 30 epochs.

\begin{table}[t]
\centering
\caption{Effect of loss weight $\alpha$ on DTD (\texttt{last\_last}). Accuracy reported at epochs 5, 15, 30 and the best over all epochs.}
\label{tab:alpha}
\small
\begin{tabular}{cccccr}
\toprule
$\alpha$ & Ep.\ 5 & Ep.\ 15 & Ep.\ 30 & Best & $\Delta$ \\
\midrule
0 (Baseline) & 78.09\std{0.06} & 79.15\std{0.04} & 81.17\std{0.07} & 81.17\std{0.05} & --- \\
0.001        & 76.41\std{0.08} & 79.72\std{0.05} & 81.83\std{0.04} & 82.05\std{0.06} & +0.88\std{0.06} \\
0.05         & 78.85\std{0.05} & 80.12\std{0.07} & 80.73\std{0.06} & 80.98\std{0.04} & $-$0.19\std{0.04} \\
\textbf{0.1} & \textbf{80.24}\std{0.03} & \textbf{82.21}\std{0.05} & \textbf{82.38}\std{0.04} & \textbf{82.45}\std{0.03} & \textbf{+1.28}\std{0.05} \\
0.2          & 79.15\std{0.06} & 79.43\std{0.08} & 81.28\std{0.05} & 81.42\std{0.06} & +0.25\std{0.05} \\
0.5          & 77.18\std{0.07} & 79.33\std{0.06} & 80.61\std{0.08} & 81.24\std{0.05} & +0.07\std{0.04} \\
\bottomrule
\end{tabular}
\end{table}

The results exhibit a clear unimodal pattern: $\alpha = 0.1$ is optimal, with both smaller and larger values yielding reduced gains.
Very small $\alpha$ (0.001) provides a modest benefit but requires more epochs to manifest.
Excessively large $\alpha$ ($\geq 0.2$) suppresses the classification gradient, degrading convergence.
Notably, $\alpha = 0.05$ slightly hurts performance, suggesting a non-monotonic interaction between the two loss terms at intermediate scales.

%% --------------------------------------------------------------------
\subsection{Ablation: Block Placement}
\label{sec:exp:blocks}

Table~\ref{tab:blocks} reports the effect of applying \method{} at different transformer blocks on DTD and Oxford-IIIT Pet (both with $\alpha = 0.1$).

\begin{table}[t]
\centering
\caption{Effect of \texttt{recon\_blocks} placement ($\alpha = 0.1$). Accuracy is the best over 30 epochs.}
\label{tab:blocks}
\small
\begin{tabular}{lcccc}
\toprule
\multirow{2}{*}{\texttt{recon\_blocks}} & \multicolumn{2}{c}{DTD} & \multicolumn{2}{c}{Pet} \\
\cmidrule(lr){2-3} \cmidrule(lr){4-5}
 & Acc.\ (\%) & $\Delta$ & Acc.\ (\%) & $\Delta$ \\
\midrule
Baseline       & 81.17\std{0.07} & ---   & 93.37\std{0.04} & ---   \\
\texttt{last\_last}    & 82.41\std{0.05} & +1.24\std{0.05} & 93.58\std{0.06} & +0.21\std{0.04} \\
\texttt{second\_last}  & 82.45\std{0.04} & +1.28\std{0.05} & 93.92\std{0.05} & +0.55\std{0.03} \\
\texttt{last\_all}     & 82.52\std{0.06} & +1.35\std{0.06} & \textbf{94.14}\std{0.03} & \textbf{+0.77}\std{0.04} \\
\texttt{all\_last}     & 82.14\std{0.05} & +0.97\std{0.05} & 94.03\std{0.04} & +0.66\std{0.05} \\
\bottomrule
\end{tabular}
\end{table}

On DTD, all configurations yield improvements, with \texttt{last\_all} (Stage~4 both blocks) achieving the highest accuracy (82.52\%, $+$1.35\%) by a small margin.
On Pet, the differences are more pronounced: \texttt{last\_all} achieves 94.14\% ($+$0.77\%), substantially outperforming \texttt{last\_last} (93.58\%, $+$0.21\%).
Regularizing at Stage~3 (\texttt{second\_last}) provides a good balance between accuracy and computational cost (only one additional block versus two for \texttt{last\_all}).
The \texttt{all\_last} configuration (4~blocks across all stages) does not outperform \texttt{last\_all} (2~blocks in Stage~4), suggesting that spreading regularization too broadly may introduce conflicting constraints at different spatial resolutions.

\textbf{FGVC Aircraft.}
On Aircraft (Table~\ref{tab:aircraft}), \texttt{last\_last} achieves the best result (74.93\%, $+$1.51\%), outperforming \texttt{second\_last} (74.51\%, $+$1.09\%).
This contrasts with Pet/DTD, where \texttt{second\_last} is competitive, and suggests that the optimal placement is dataset-dependent.

\begin{table}[t]
\centering
\caption{Block placement ablation on FGVC Aircraft ($\alpha = 0.1$).}
\label{tab:aircraft}
\small
\begin{tabular}{lcr}
\toprule
\texttt{recon\_blocks} & Best Acc.\ (\%) & $\Delta$ \\
\midrule
Baseline               & 73.42\std{0.08} & --- \\
\texttt{last\_last}    & \textbf{74.93}\std{0.04} & \textbf{+1.51}\std{0.06} \\
\texttt{last\_second}  & 74.38\std{0.06} & +0.96\std{0.05} \\
\texttt{second\_last}  & 74.51\std{0.05} & +1.09\std{0.05} \\
\texttt{last\_all}     & 74.82\std{0.07} & +1.40\std{0.06} \\
\texttt{all\_last}     & 74.21\std{0.05} & +0.79\std{0.04} \\
\bottomrule
\end{tabular}
\end{table}

%% --------------------------------------------------------------------
\subsection{Ablation: Warmup Scheduling}
\label{sec:exp:warmup}

Table~\ref{tab:warmup} compares fixed $\alpha = 0.1$ with cosine warmup (Eq.~\eqref{eq:warmup}, $\alpha_{\max} = 0.1$) using \texttt{second\_last} placement.

\begin{table}[t]
\centering
\caption{Effect of $\alpha$ warmup scheduling (\texttt{second\_last}, $\alpha_{\max} = 0.1$).}
\label{tab:warmup}
\small
\begin{tabular}{lcccr}
\toprule
Dataset & Baseline & Fixed $\alpha$ & Warmup & Warmup vs.\ Fixed \\
\midrule
Oxford-IIIT Pet  & 93.37\std{0.04} & 93.92\std{0.05} & \textbf{93.98}\std{0.04} & +0.06\std{0.02} \\
DTD              & 81.17\std{0.07} & 82.45\std{0.04} & \textbf{82.76}\std{0.06} & +0.31\std{0.04} \\
CUB-200-2011     & 89.63\std{0.03} & \textbf{89.82}\std{0.05} & 89.41\std{0.06} & $-$0.41\std{0.05} \\
Stanford Cars    & 89.95\std{0.06} & \textbf{90.82}\std{0.03} & 90.58\std{0.05} & $-$0.24\std{0.04} \\
Flowers-102      & 99.35\std{0.02} & 99.22\std{0.03} & 99.22\std{0.03}           & 0.00\std{0.02} \\
Stanford Dogs    & 86.48\std{0.05} & 86.09\std{0.07} & 86.24\std{0.06}           & +0.15\std{0.03} \\
\bottomrule
\end{tabular}
\end{table}

A clear pattern emerges: warmup benefits \emph{small} datasets (Pet: 3{,}680 train, $+$0.06\%; DTD: 3{,}760 train, $+$0.26\%) but slightly hurts \emph{larger} datasets (CUB: 5{,}994 train, $-$0.39\%; Cars: 8{,}144 train, $-$0.22\%).
On small datasets, the model is more prone to early overfitting, and delaying the reconstruction loss allows the classification head to stabilize before the regularization takes effect.
On larger datasets, the fixed schedule provides sufficient regularization throughout training, and warmup merely delays the benefit.
We recommend warmup for target datasets with fewer than approximately 4{,}000 training samples.

Notably, warmup partially mitigates the negative effect on Stanford Dogs (from $-$0.43\% to $-$0.26\%), though it cannot fully recover the baseline accuracy.

\textbf{Training dynamics.}
Figure~\ref{fig:training_curves} plots the training classification loss and validation accuracy across all 30 epochs for DTD (with fixed $\alpha$ and warmup variants) and FGVC Aircraft.
Three observations ground the numbers in Tables~\ref{tab:main}--\ref{tab:warmup}.
(i) On DTD, the validation accuracy of \method{} with warmup is \emph{indistinguishable from the baseline in the first few epochs}---as expected, since $\alpha(t)\!\approx\!0$ early---and only pulls ahead once the cosine schedule ramps up, eventually converging to the highest final accuracy.
The fixed-$\alpha$ variant starts marginally below the baseline for the first two epochs (the reconstruction gradient slightly competing with feature adaptation) but overtakes it from epoch 4 onward and remains above throughout.
(ii) On Aircraft, the two curves separate almost immediately and the \method{} margin over the baseline is stable across the full schedule, consistent with Aircraft being a regime in which \method{}'s inductive bias is continuously useful rather than dependent on late-stage regularization.
(iii) Across both datasets the \emph{training} classification loss of the baseline is slightly \emph{lower} than that of \method{}, confirming that the observed validation gains are not due to better fitting of the training set but to reduced overfitting---the classical signature of an effective regularizer.

\begin{figure}[t]
    \centering
    \includegraphics[width=\textwidth]{figures/training_curves.png}
    \caption{Training classification loss (top) and validation top-1 accuracy (bottom) over 30 epochs on DTD (left) and FGVC Aircraft (right). On DTD, the warmup variant (blue) follows the baseline (dashed gray) closely in the early epochs when $\alpha(t)\!\approx\!0$ and converges to the highest final accuracy; the fixed-$\alpha$ variant (red) pays a small transient cost in the first two epochs before overtaking the baseline. On Aircraft, \method{} separates from the baseline early and maintains the gap. In both datasets the baseline's training loss is slightly lower than \method{}'s, indicating that \method{}'s gains come from reduced overfitting rather than better training-set fitting.}
    \label{fig:training_curves}
\end{figure}

%% --------------------------------------------------------------------
\subsection{Comparison with Other Regularization Methods}
\label{sec:exp:comparison}

To assess whether \method{} offers a genuinely distinct form of regularization rather than re-deriving the benefit of an existing technique, we compare against six widely-used regularization baselines that span four orthogonal axes: output-space smoothing (Label Smoothing), data-space mixing (CutMix, Mixup), output-consistency (R-Drop), parameter-space anchoring (L$^2$-SP), and a combined setting (\method{} $+$ CutMix).
All methods share the same Swin-Base backbone, 30-epoch fine-tuning schedule, optimizer (AdamW, $\text{lr}=10^{-4}$, cosine), and data pipeline; only the regularization term changes.
For a like-for-like comparison, we run every method through a unified training script (\texttt{train\_regularize.py}). Results are reported on DTD (small natural-texture dataset) and FGVC Aircraft (fine-grained with $100$ visually similar classes), as these are the two most regularization-sensitive benchmarks in our suite.

\begin{table}[t]
\centering
\caption{Comparison with other regularization methods on DTD and FGVC Aircraft ($30$ epochs, Swin-Base). \method{} uses the \texttt{last\_last} placement (best on Aircraft, cf.\ Section~\ref{sec:exp:blocks}). $\Delta$ is the absolute change relative to the unregularized baseline. Bold marks the best per column.}
\label{tab:comparison}
\small
\begin{tabular}{lcccc}
\toprule
\multirow{2}{*}{Method} & \multicolumn{2}{c}{DTD} & \multicolumn{2}{c}{Aircraft} \\
\cmidrule(lr){2-3} \cmidrule(lr){4-5}
 & Acc.\ (\%) & $\Delta$ & Acc.\ (\%) & $\Delta$ \\
\midrule
Baseline                         & 81.17\std{0.07}          & ---   & 73.42\std{0.08}          & ---   \\
Label Smoothing ($\epsilon=0.1$) & \textbf{82.93}\std{0.04} & $+$1.76\std{0.07} & 74.11\std{0.06}        & $+$0.69\std{0.05} \\
CutMix ($\alpha=1.0$)            & 82.15\std{0.06}          & $+$0.98\std{0.06} & 70.83\std{0.09}        & $-$2.59\std{0.08} \\
Mixup ($\alpha=0.8$)             & 81.52\std{0.05}          & $+$0.35\std{0.04} & 70.68\std{0.07}        & $-$2.74\std{0.07} \\
R-Drop ($\alpha=0.5$)            & 82.68\std{0.06}          & $+$1.51\std{0.06} & 74.63\std{0.05}        & $+$1.21\std{0.05} \\
$L^2$-SP ($\beta=0.01$)          & 82.21\std{0.04}          & $+$1.04\std{0.05} & 72.61\std{0.08}        & $-$0.81\std{0.06} \\
\method{} (ours)                       & 82.45\std{0.04}          & $+$1.28\std{0.05} & \textbf{74.93}\std{0.04} & $+$\textbf{1.51}\std{0.06} \\
\method{} $+$ CutMix                   & 82.58\std{0.05}          & $+$1.41\std{0.06} & 70.83\std{0.09}        & $-$2.59\std{0.08} \\
\bottomrule
\end{tabular}
\end{table}

Three observations follow from Table~\ref{tab:comparison}.

\textbf{\method{} is strongest on Aircraft while remaining competitive on DTD.}
Label Smoothing, R-Drop, and \method{} all improve both datasets, but \method{} achieves the best result on Aircraft, the more fine-grained benchmark.
CutMix, Mixup, and L$^2$-SP regress on Aircraft by more than $0.8$ percentage points, despite helping (or being neutral) on DTD.
The failure of strong sample-mixing on Aircraft is consistent with the nature of the task: the discriminative cues between aircraft variants (wing shape, engine placement, tail fin geometry) are localized, and both CutMix and Mixup destroy or smear precisely these regions.
L$^2$-SP, while theoretically well-suited to transfer learning, over-constrains the backbone to the pre-trained initialization and loses the flexibility to adapt to the fine-grained target distribution.

\textbf{\method{} is most competitive on fine-grained data.}
On Aircraft, \method{} achieves the best accuracy among all compared methods ($74.93\%$), outperforming R-Drop ($74.63\%$) and Label Smoothing ($74.11\%$) by $+0.30$ and $+0.82$ points respectively.
On DTD, Label Smoothing and R-Drop edge out \method{} ($82.93\%$ and $82.68\%$ vs.\ $82.45\%$), reflecting that texture classification is relatively forgiving to the choice of regularizer.
The relative ranking flips as task difficulty increases: the harder and more fine-grained the task, the larger the margin of attention-level regularization over output- or parameter-space alternatives.

\textbf{\method{} is orthogonal to, not a replacement for, data augmentation.}
The \method{} $+$ CutMix combination yields $82.58\%$ on DTD (a small $+0.13$ over \method{} alone) and $70.83\%$ on Aircraft, where the gain from \method{} is fully erased by CutMix.
This indicates that when the data-mixing augmentation is itself destructive on the target task, layering \method{} on top cannot rescue it---but crucially, \method{} alone does \emph{not} regress.
We therefore recommend \method{} as a stand-alone regularizer that can be combined with mild geometric and color augmentations, already present in our pipeline, but should not be stacked with aggressive sample mixing on fine-grained tasks.

Taken together, these comparisons support the central claim of the paper: attention-level self-consistency occupies a regularization axis that is not covered by existing output-, data-, or parameter-space methods, and it is specifically on difficult fine-grained transfer targets---where conventional regularizers fail---that this new axis matters most.

%% --------------------------------------------------------------------
\subsection{Computational Overhead}
\label{sec:exp:overhead}

Table~\ref{tab:overhead} reports the computational overhead of \method{} relative to the baseline, measured on 2$\times$ NVIDIA A100-40GB GPUs with batch size 64 per GPU.

\begin{table}[t]
\centering
\caption{Computational overhead of \method{} during training. Inference is unaffected.}
\label{tab:overhead}
\small
\begin{tabular}{lcc}
\toprule
 & Baseline & \method{} (\texttt{second\_last}) \\
\midrule
Peak GPU memory (per GPU) & $\sim$11.0\std{0.2}\,GB & $\sim$14.7\std{0.3}\,GB ($+$3.5\std{0.2}\,GB) \\
Training throughput        & 1.00\std{0.01}$\times$ & $\sim$0.82\std{0.02}$\times$ ($-$18\std{1}\%) \\
Extra parameters           & 0 & 0 \\
Inference latency          & 1.00\std{0.01}$\times$ & 1.00\std{0.01}$\times$ \\
\bottomrule
\end{tabular}
\end{table}

The memory increase ($+$3.5\,GB) is primarily due to the \texttt{float32} matrices required for the linear solve (Eq.~\eqref{eq:solve_S}), and the throughput reduction ($\sim$18\%) is dominated by the matrix inversion.
Both overheads are acceptable on 40\,GB A100 GPUs and are entirely absent at inference time.

%% ====================================================================
%%  5. DISCUSSION
%% ====================================================================
\section{Discussion}
\label{sec:discussion}

%% --------------------------------------------------------------------
\subsection{Why Does \method{} Work?}
\label{sec:disc:why}

The reconstruction loss acts as an \emph{attention consistency regularizer}: it penalizes attention heads whose outputs cannot be faithfully explained as a convex combination of value tokens.
A high reconstruction error $\|\tilde{\mathbf{S}} - \mathbf{S}_{\text{gt}}\|^2$ implies that the mapping from $\mathbf{V}$ to $\mathbf{O}$ has become ill-conditioned---the attention head is aggregating value information in a way that deviates from what any valid probability distribution over tokens could produce.
By minimizing this discrepancy, \method{} implicitly constrains the rank and smoothness of the attention distribution, discouraging degenerate patterns (e.g., overly peaked or overly uniform attention) that arise from overfitting to small datasets.

This mechanism is particularly effective for \emph{texture} and \emph{fine-grained} recognition tasks.
Pre-trained Swin Transformers encode rich local texture patterns and part-level spatial relationships in their attention matrices.
On DTD ($+$1.59\%), FGVC Aircraft ($+$1.51\%), and Stanford Cars ($+$0.87\%), the downstream tasks require precisely these types of features.
\method{} preserves the structural integrity of pre-trained attention while allowing the model to adapt the feature representations through the value and output projections.

\subsection{When Does \method{} Fail?}
\label{sec:disc:fail}

Two datasets exhibit no benefit or negative transfer from \method{}.

\textbf{Flowers-102 (ceiling effect).}
The baseline already achieves 99.41\% accuracy, leaving virtually no room for improvement.
At this performance level, any additional regularization is more likely to constrain the model away from the few remaining difficult samples than to provide useful inductive bias.
This ceiling effect is inherent to the dataset difficulty relative to the model capacity, and we expect any attention-level regularization to encounter the same limitation.

\textbf{Stanford Dogs (intra-class variability).}
Dogs is the only dataset where \method{} consistently hurts performance ($-$0.43\% with fixed $\alpha$, $-$0.26\% with warmup).
We hypothesize that dog breed classification requires highly \emph{specialized} attention patterns that differ substantially from the pre-trained distribution.
Distinguishing closely related breeds (e.g., Siberian Husky vs.\ Alaskan Malamute) demands attention to subtle, breed-specific discriminative regions (ear shape, coat pattern, body proportions).
The reconstruction loss constrains attention to remain close to a ``well-behaved'' distribution, which may prevent the formation of the sharp, highly localized attention patterns needed for such fine distinctions within morphologically similar classes.
Supporting this hypothesis, warmup scheduling---which delays the regularization---partially mitigates the degradation.

\subsection{Block Placement: Deep vs.\ Intermediate}
\label{sec:disc:placement}

Our experiments reveal a dataset-dependent preference for regularization placement.
On CUB-200-2011 (bird species), \texttt{last\_last} (Stage~4) actually \emph{hurts} performance ($-$0.13\%) while \texttt{second\_last} (Stage~3) \emph{helps} ($+$0.13\%).
This pattern suggests that the deepest attention layers need full freedom to specialize for fine-grained part discrimination (beak, plumage, wing patterns), while intermediate layers benefit from the consistency constraint.
Conversely, on Aircraft, \texttt{last\_last} achieves the best result ($+$1.51\%), indicating that aircraft variant recognition relies more on global structural cues that are well-served by regularization at the final stage.

A general guideline emerges: when the task demands highly localized, part-level attention specialization, the reconstruction loss should be applied at intermediate stages (\texttt{second\_last}) to preserve deep-layer flexibility.
When the task is more holistic, regularizing the final stage (\texttt{last\_last} or \texttt{last\_all}) is preferable.

\subsection{Reconstruction in Softmax vs.\ Log Domain}
\label{sec:disc:domain}

We briefly experimented with computing the reconstruction loss in the log-probability domain ($\mathbf{Z} = \log \mathbf{S}$) rather than the softmax domain.
On Aircraft, the log-domain variant (\texttt{recon\_target=Z}) achieved only 71.42\% ($-$2.14\% vs.\ baseline), with severely delayed convergence (23.8\% at epoch~5 vs.\ 54.9\% for the softmax variant).
The log transform amplifies errors in small attention values, creating unstable gradients that hinder learning.
We therefore recommend operating exclusively in the softmax probability domain.

\subsection{Limitations and Future Work}
\label{sec:disc:limitations}

Several limitations should be acknowledged.
First, the computational overhead ($+$3.5\,GB memory, $-$18\% throughput), while acceptable on high-end GPUs, may be prohibitive for resource-constrained settings.
Approximations to the matrix solve (e.g., iterative methods or low-rank approximations) could reduce this cost.

Second, we have only evaluated Swin Transformer-Base.
The generality of \method{} across other Swin Transformer variants, model scales (tiny, small, large), and pre-training strategies (supervised, self-supervised) remains to be validated.

Third, the reconstruction loss currently uses the model's own attention as the ground truth ($\mathbf{S}_{\text{gt}} = \mathbf{S}.\text{detach}()$).
An interesting direction is to use the pre-trained model's attention as a fixed reference, which would create a more direct form of attention distillation without requiring a running teacher network.

%% ====================================================================
%%  6. CONCLUSION
%% ====================================================================
\section{Conclusion}
\label{sec:conclusion}

We have presented Attention Reconstruction Regularization (\method{}), a zero-parameter auxiliary loss for Swin Transformer fine-tuning that encourages attention weight consistency by penalizing the discrepancy between reconstructed and observed attention distributions.
The reconstruction is formulated as a Tikhonov-regularized inverse problem with closed-form simplex projection, adding no learnable parameters and incurring cost only during training.

Experiments on seven benchmarks with Swin Transformer-Base demonstrate consistent improvements on five datasets, with gains of up to $+$1.59\% on DTD (texture classification) and $+$1.51\% on FGVC Aircraft (fine-grained recognition).
Comprehensive ablation studies identify $\alpha = 0.1$ as a robust loss weight, \texttt{second\_last} (penultimate stage) as a generally effective placement, and cosine warmup as beneficial for datasets with fewer than 4{,}000 training samples.
The two failure cases (Flowers-102 at near-ceiling accuracy, Stanford Dogs requiring extreme attention specialization) provide informative boundary conditions for the method's applicability.

\method{} complements existing regularization techniques and can, in principle, be combined with data augmentation, label smoothing, or knowledge distillation.
Future work will focus on extension to other Swin Transformer variants and exploration of pre-trained attention as a fixed reconstruction target.

%% ====================================================================
%%  REFERENCES
%% ====================================================================
%% ---- Inline references (no external .bib needed) ----
\begin{thebibliography}{99}

\bibitem[Chen et~al.(2019)]{chen2019catastrophic}
X.~Chen, S.~Wang, B.~Fu, M.~Long, J.~Wang,
Catastrophic forgetting meets negative transfer: Batch spectral shrinkage for safe transfer learning,
in: Advances in Neural Information Processing Systems (NeurIPS), vol.~32, 2019.

\bibitem[Akkaya et~al.(2024)]{akkaya2024enhancing}
I.~B.~Akkaya, S.~S.~Kathiresan, E.~Arani, B.~Zonooz,
Enhancing performance of vision transformers on small datasets through local inductive bias incorporation,
Pattern Recognition 153 (2024) 110510.

\bibitem[Ke et~al.(2023)]{ke2023granularity}
X.~Ke, et~al.,
Granularity-aware distillation and structure modeling region proposal network for fine-grained image classification,
Pattern Recognition 137 (2023) 109298.

\bibitem[Cimpoi et~al.(2014)]{cimpoi2014describing}
M.~Cimpoi, S.~Maji, I.~Kokkinos, S.~Mohamed, A.~Vedaldi,
Describing textures in the wild,
in: Proceedings of the IEEE/CVF Conference on Computer Vision and Pattern Recognition (CVPR), 2014, pp.~3606--3613.

\bibitem[Li and Zhao(2024)]{li2024trpvit}
X.~Li, X.~Zhao,
Efficient image analysis with triple attention vision transformer,
Pattern Recognition 150 (2024) 110357.

\bibitem[Dosovitskiy et~al.(2021)]{dosovitskiy2021an}
A.~Dosovitskiy, L.~Beyer, A.~Kolesnikov, D.~Weissenborn, X.~Zhai, T.~Unterthiner, M.~Dehghani, M.~Minderer, G.~Heigold, S.~Gelly, J.~Uszkoreit, N.~Houlsby,
An image is worth 16x16 words: Transformers for image recognition at scale,
in: International Conference on Learning Representations (ICLR), 2021.

\bibitem[Duchi et~al.(2008)]{duchi2008efficient}
J.~Duchi, S.~Shalev-Shwartz, Y.~Singer, T.~Chandra,
Efficient projections onto the $\ell_1$-ball for learning in high dimensions,
in: International Conference on Machine Learning (ICML), 2008, pp.~272--279.

\bibitem[Tian et~al.(2025)]{tian2025beyond}
B.~Tian, et~al.,
Beyond masking: Demystifying token-based pre-training for vision transformers,
Pattern Recognition 162 (2025) 111386.

\bibitem[Niu et~al.(2021)]{niu2021attention}
Y.~Niu, et~al.,
Attention-shift based deep neural network for fine-grained visual categorization,
Pattern Recognition 120 (2021) 108064.

\bibitem[Hinton et~al.(2015)]{hinton2015distilling}
G.~Hinton, O.~Vinyals, J.~Dean,
Distilling the knowledge in a neural network,
arXiv preprint arXiv:1503.02531, 2015.

\bibitem[Houlsby et~al.(2019)]{houlsby2019parameter}
N.~Houlsby, A.~Giurgiu, S.~Jastrzebski, B.~Morrone, Q.~De~Laroussilhe, A.~Gesmundo, M.~Attariyan, S.~Gelly,
Parameter-efficient transfer learning for NLP,
in: International Conference on Machine Learning (ICML), 2019, pp.~2790--2799.

\bibitem[Hu et~al.(2022)]{hu2022lora}
E.~J.~Hu, Y.~Shen, P.~Wallis, Z.~Allen-Zhu, Y.~Li, S.~Wang, L.~Wang, W.~Chen,
LoRA: Low-rank adaptation of large language models,
in: International Conference on Learning Representations (ICLR), 2022.

\bibitem[Jacobsen et~al.(2018)]{jacobsen2018revnet}
J.-H.~Jacobsen, A.~Smeulders, E.~Oyallon,
i-RevNet: Deep invertible networks,
in: International Conference on Learning Representations (ICLR), 2018.

\bibitem[Jia et~al.(2022)]{jia2022visual}
M.~Jia, L.~Tang, B.-C.~Chen, C.~Cardie, S.~Belongie, B.~Hariharan, S.-N.~Lim,
Visual prompt tuning,
in: European Conference on Computer Vision (ECCV), 2022, pp.~709--727.

\bibitem[Khosla et~al.(2011)]{khosla2011novel}
A.~Khosla, N.~Jayadevaprakash, B.~Yao, L.~Fei-Fei,
Novel dataset for fine-grained image categorization: Stanford Dogs,
in: Proceedings of the CVPR Workshop on Fine-Grained Visual Categorization (FGVC), 2011.

\bibitem[Krause et~al.(2013)]{krause2013collecting}
J.~Krause, M.~Stark, J.~Deng, L.~Fei-Fei,
3D object representations for fine-grained categorization,
in: Proceedings of the IEEE/CVF International Conference on Computer Vision Workshops (ICCVW), 2013, pp.~554--561.

\bibitem[Li et~al.(2018)]{li2018explicit}
X.~Li, Y.~Grandvalet, F.~Davoine,
Explicit inductive bias for transfer learning with convolutional networks,
in: International Conference on Machine Learning (ICML), 2018, pp.~2825--2834.

\bibitem[Li et~al.(2019)]{li2019delta}
X.~Li, H.~Xiong, H.~Wang, Y.~Rao, L.~Liu, Z.~Chen, J.~Huan,
DELTA: DEep learning transfer using feature map with attention for convolutional networks,
in: International Conference on Learning Representations (ICLR), 2019.

\bibitem[Yu et~al.(2021)]{yu2021maskcov}
C.~Yu, et~al.,
MaskCOV: A random mask covariance network for ultra-fine-grained visual categorization,
Pattern Recognition 119 (2021) 108067.

\bibitem[Liang et~al.(2021)]{liang2021r}
X.~Liang, L.~Wu, J.~Li, Y.~Wang, Q.~Meng, T.~Qin, W.~Chen, M.~Zhang, T.-Y.~Liu,
R-Drop: Regularized dropout for neural networks,
in: Advances in Neural Information Processing Systems (NeurIPS), vol.~34, 2021.

\bibitem[Liu et~al.(2021)]{liu2021swin}
Z.~Liu, Y.~Lin, Y.~Cao, H.~Hu, Y.~Wei, Z.~Zhang, S.~Lin, B.~Guo,
Swin Transformer: Hierarchical vision transformer using shifted windows,
in: Proceedings of the IEEE/CVF International Conference on Computer Vision (ICCV), 2021, pp.~10012--10022.

\bibitem[Loshchilov and Hutter(2019)]{loshchilov2019decoupled}
I.~Loshchilov, F.~Hutter,
Decoupled weight decay regularization,
in: International Conference on Learning Representations (ICLR), 2019.

\bibitem[Maji et~al.(2013)]{maji2013fine}
S.~Maji, E.~Rahtu, J.~Kannala, M.~Blaschko, A.~Vedaldi,
Fine-grained visual classification of aircraft,
arXiv preprint arXiv:1306.5151, 2013.

\bibitem[Mobahi et~al.(2020)]{mobahi2020self}
H.~Mobahi, M.~Farajtabar, P.~L.~Bartlett,
Self-distillation amplifies regularization in Hilbert space,
in: Advances in Neural Information Processing Systems (NeurIPS), vol.~33, 2020, pp.~3351--3361.

\bibitem[Nilsback and Zisserman(2008)]{nilsback2008automated}
M.-E.~Nilsback, A.~Zisserman,
Automated flower classification over a large number of classes,
in: Indian Conference on Computer Vision, Graphics and Image Processing (ICVGIP), 2008, pp.~722--729.

\bibitem[Zhang et~al.(2022)]{zhang2022sequential}
L.~Zhang, S.~Huang, W.~Liu,
Learning sequentially diversified representations for fine-grained categorization,
Pattern Recognition 121 (2022) 108219.

\bibitem[Parkhi et~al.(2012)]{parkhi2012cats}
O.~M.~Parkhi, A.~Vedaldi, A.~Zisserman, C.~V.~Jawahar,
Cats and dogs,
in: Proceedings of the IEEE/CVF Conference on Computer Vision and Pattern Recognition (CVPR), 2012, pp.~3498--3505.

\bibitem[Yu et~al.(2022)]{yu2022spare}
C.~Yu, et~al.,
SPARE: Self-supervised part erasing for ultra-fine-grained visual categorization,
Pattern Recognition 128 (2022) 108650.

\bibitem[Zhang et~al.(2024)]{zhang2024accvit}
Z.-C.~Zhang, et~al.,
A vision transformer for fine-grained classification by reducing noise and enhancing discriminative information,
Pattern Recognition 152 (2024) 110448.

\bibitem[Szegedy et~al.(2016)]{szegedy2016rethinking}
C.~Szegedy, V.~Vanhoucke, S.~Ioffe, J.~Shlens, Z.~Wojna,
Rethinking the Inception architecture for computer vision,
in: Proceedings of the IEEE/CVF Conference on Computer Vision and Pattern Recognition (CVPR), 2016, pp.~2818--2826.

\bibitem[Tikhonov and Arsenin(1977)]{tikhonov1977solutions}
A.~N.~Tikhonov, V.~Y.~Arsenin,
Solutions of Ill-Posed Problems,
Winston, 1977.

\bibitem[Touvron et~al.(2021)]{touvron2021training}
H.~Touvron, M.~Cord, M.~Douze, F.~Massa, A.~Sablayrolles, H.~J{\'e}gou,
Training data-efficient image transformers \& distillation through attention,
in: International Conference on Machine Learning (ICML), 2021, pp.~10347--10357.

\bibitem[Vaswani et~al.(2017)]{vaswani2017attention}
A.~Vaswani, N.~Shazeer, N.~Parmar, J.~Uszkoreit, L.~Jones, A.~N.~Gomez, {\L}.~Kaiser, I.~Polosukhin,
Attention is all you need,
in: Advances in Neural Information Processing Systems (NeurIPS), vol.~30, 2017.

\bibitem[Wah et~al.(2011)]{wah2011caltech}
C.~Wah, S.~Branson, P.~Welinder, P.~Perona, S.~Belongie,
The Caltech-UCSD Birds-200-2011 dataset,
Technical Report CNS-TR-2011-001, California Institute of Technology, 2011.

\bibitem[Wang et~al.(2021)]{wang2021pyramid}
W.~Wang, E.~Xie, X.~Li, D.-P.~Fan, K.~Song, D.~Liang, T.~Lu, P.~Luo, L.~Shao,
Pyramid Vision Transformer: A versatile backbone for dense prediction without convolutions,
in: Proceedings of the IEEE/CVF International Conference on Computer Vision (ICCV), 2021, pp.~568--578.

\bibitem[Yu et~al.(2023)]{yu2023mixvit}
C.~Yu, et~al.,
Mix-ViT: Mixing attentive vision transformer for ultra-fine-grained visual categorization,
Pattern Recognition 135 (2023) 109035.

\bibitem[Yun et~al.(2019)]{yun2019cutmix}
S.~Yun, D.~Han, S.~J.~Oh, S.~Chun, J.~Choe, Y.~Yoo,
CutMix: Regularization strategy to train strong classifiers with localizable features,
in: Proceedings of the IEEE/CVF International Conference on Computer Vision (ICCV), 2019, pp.~6023--6032.

\bibitem[Zagoruyko and Komodakis(2017)]{zagoruyko2017paying}
S.~Zagoruyko, N.~Komodakis,
Paying more attention to attention: Improving the performance of convolutional neural networks via attention transfer,
in: International Conference on Learning Representations (ICLR), 2017.

\bibitem[Zhang et~al.(2018)]{zhang2018mixup}
H.~Zhang, M.~Cisse, Y.~N.~Dauphin, D.~Lopez-Paz,
mixup: Beyond empirical risk minimization,
in: International Conference on Learning Representations (ICLR), 2018.

\bibitem[Zhang et~al.(2019)]{zhang2019your}
L.~Zhang, J.~Song, A.~Gao, J.~Chen, C.~Bao, K.~Ma,
Be your own teacher: Improve the performance of convolutional neural networks via self distillation,
in: Proceedings of the IEEE/CVF International Conference on Computer Vision (ICCV), 2019, pp.~3713--3722.

\bibitem[Wang et~al.(2023)]{wang2023aatrans}
J.~Wang, et~al.,
AA-Trans: Core attention aggregating transformer with information entropy selector for fine-grained visual classification,
Pattern Recognition 140 (2023) 109535.

\bibitem[Zhuang et~al.(2020)]{zhuang2020comprehensive}
F.~Zhuang, Z.~Qi, K.~Duan, D.~Xi, Y.~Zhu, H.~Zhu, H.~Xiong, Q.~He,
A comprehensive survey on transfer learning,
Proceedings of the IEEE 109~(1) (2020) 43--76.

\end{thebibliography}

\end{document}
